Assume \(\mathcal R=\varnothing\), \(\mathcal V^{\dagger}=\varnothing\), and the label set is empty.  There is then exactly one empty context, no selection transition, and no nontrivial label license.  The observed record is the substantive record.  Unfolding \(\text{def:trace-state}\) and \(\text{def:tcf-procedure}\), the only graphical transitions are ancestral marginalization, district decomposition, and uncolored fixing.  These are exactly the transitions of ordinary \(\operatorname{ID}\), and the extra context and color components remain vacuous.  Hence every rational graphical trace returned by TCF is syntactically the corresponding ordinary fixing expression.

Suppose ordinary \(\operatorname{ID}\) returns \(f_a^{\mathrm{ID}}(P_{\mathrm{obs}})(y)\).  By the cited soundness theorem \(\text{lem:ordinary-id-soundness}\),
\[
Q_a^{\mathcal M}(y)=f_a^{\mathrm{ID}}(P_{\mathrm{obs}})(y)
\quad\text{for every }\mathcal M\in\mathfrak F(\mathcal G,P_{\mathrm{obs}}).
\tag{1}
\]
The right side is fixed by the observed law, so the causal coordinate is constant on the fiber.  By \(\text{lem:semialgebraic-fiber-dichotomy}\), the unique TCF cell is identifying.  The ordinary trace is among the enumerated candidates, and its equality with the unique value is valid on every compatible point of the cell by (1); therefore the whole-cell annotation test in \(\text{def:tcf-procedure}\) accepts it.  Applying \(\text{thm:tcf-soundness}\) gives
\[
f_{a,j}(P_{\mathrm{obs}})(y)
=Q_a^{\mathcal M}(y)
=f_a^{\mathrm{ID}}(P_{\mathrm{obs}})(y).
\tag{2}
\]
The semialgebraic test is pointwise in an observed law, whereas the ordinary ID graphwise dichotomy is uniform over compatible laws.  Equations (1)--(2) prove the asserted one-way reduction and do not entail the excluded converse.